%% toc: P2 23 23.5,23.6,23.7,23.8
\section{Full second-order validity has no recursive enumeration: a direct proof}
\label{sec:p02-w02-p02-041-direct}

\emph{Purpose.}  This section proves, by a reduction that stays inside
second-order logic, that the set of full second-order validities over the
arithmetic signature is not the domain of any partial recursive function,
and hence that no calculus with a recursively enumerable set of theorems is
sound and complete for full validity.  The proof does not pass through the
arithmetical hierarchy, through \cref{thm:p04-ce-sigma-one}, or through
true arithmetic; it describes the graph of a partial recursive function
directly by second-order conditions whose only solution, in the standard
model, is that graph.

\emph{Relation to the surrounding material.}  The typed higher-order calculus
of this chapter, with its primitive constants, conversion rules, axioms and
rules (comprehension by abstraction, choice and descriptions by a selection
constant, extensionality by an axiom schema), is Church's system of
\cref{sec:p02-w02-p02-040-church}; its general models and their soundness and
completeness through a first-order encoding and decoding are
\cref{thm:p02-w02-p02-040-church-general-soundness,thm:p02-w02-p02-040-church-encoding,thm:p02-w02-p02-040-church-general-completeness};
the relational calculus with its comprehension and choice schemas is
\cref{def:p02-sol-derived-calculus}, and the completeness of its native finite
rule list is \cref{thm:p02-w02-p02-040-sol-native}.  The nonaxiomatizability
of full validity is \cref{thm:p02-w02-p02-040-no-effective-calculus}, whose
proof of clause~3 rests on \cref{thm:p04-ce-sigma-one}, a result proved there
only in outline.  \Cref{thm:p02-w02-p02-041-no-recursive-domain} below
proves clauses~4 and~5 of that theorem by a route that uses neither clause~3
nor \cref{thm:p04-ce-sigma-one}.

\emph{Inputs.}  The relational second-order language over
\(\{0,S,+,\times\}\) with full semantics
(\cref{sec:second-order-logic-01,sec:second-order-logic-02}); the sentence
\(\mathsf{PA}_2\) of \cref{def:p02-pa2} and its full categoricity
(\cref{thm:p02-w02-p02-040-pa2-categorical}); the codes of second-order
formulas of \cref{def:p02-w02-p02-040-so-coding}, built on the list and
constructor coding of \cref{sec:arithmetic-coding-04}; primitive recursive
functions (\cref{sec:turing-machines-02}) and partial recursive functions
(\cref{def:p04-partial-recursive}).  For the formulation with Turing
machines, also \cref{thm:p04-tm-partial-recursive}.
\emph{Conclusions.}
\Cref{thm:p02-w02-p02-041-coverage,thm:p02-w02-p02-041-iso,thm:p02-w02-p02-041-graphs,thm:p02-w02-p02-041-reduction,thm:p02-w02-p02-041-effective,thm:p02-w02-p02-041-no-recursive-domain}.
\emph{Standing.}  The reduction and all proofs are the handbook's.  Every
step is proved in the annex except the passage from Turing-machine
computability to partial recursiveness in
\cref{thm:p02-w02-p02-041-no-recursive-domain}(c) and~(d), which uses the
cited \cref{thm:p04-tm-partial-recursive}; clauses (a) and (b) are free of
that citation.

\subsection{Recursion terms}
\label{sec:p02-w02-p02-041-terms}

\begin{hbdefinition}{def:p02-w02-p02-041-terms}{Recursion terms, their arities and codes}
\emph{Recursion terms} \(\rho\) and their arities \(\mathrm{ar}(\rho)\in\N\)
are generated by the following six clauses.  Each clause also fixes the code
\(\ulcorner\rho\urcorner\), using the constructor
\(N_k(a_0,\ldots,a_{r-1})=C(k,[a_0,\ldots,a_{r-1}])\) of
\cref{sec:arithmetic-coding-04} with the tags \(30\) to \(35\), which occur
in no other sort used in this chapter.
\begin{enumerate}
 \item For every \(k\in\N\), \(\mathbf Z^k\) is a term of arity \(k\), with
   code \(N_{30}(k)\).
 \item \(\mathbf S\) is a term of arity \(1\), with code \(N_{31}()\).
 \item For \(1\le i\le k\), \(\mathbf P^k_i\) is a term of arity \(k\), with
   code \(N_{32}(k,i)\).
 \item If \(k,m\in\N\), \(\gamma\) has arity \(m\) and
   \(\eta_1,\ldots,\eta_m\) each have arity \(k\), then
   \(\mathbf C^k(\gamma;\eta_1,\ldots,\eta_m)\) is a term of arity \(k\), with
   code \(N_{33}(k,\ulcorner\gamma\urcorner,[\ulcorner\eta_1\urcorner,\ldots,\ulcorner\eta_m\urcorner])\).
 \item If \(\gamma\) has arity \(k\) and \(\eta\) has arity \(k+2\), then
   \(\mathbf R(\gamma,\eta)\) is a term of arity \(k+1\), with code
   \(N_{34}(\ulcorner\gamma\urcorner,\ulcorner\eta\urcorner)\).
 \item If \(\gamma\) has arity \(k+1\), then \(\mathbf M(\gamma)\) is a term
   of arity \(k\), with code \(N_{35}(\ulcorner\gamma\urcorner)\).
\end{enumerate}
The terms \(\gamma\) and \(\eta_j\), \(\eta\) displayed in clauses 4--6 are the
\emph{immediate subterms}.  The term \(\rho\) \emph{denotes} the partial map
\(f_\rho:\N^{\mathrm{ar}(\rho)}\rightharpoonup\N\) given by
\(f_{\mathbf Z^k}(\vec x)=0\), \(f_{\mathbf S}(x)=x+1\),
\(f_{\mathbf P^k_i}(\vec x)=x_i\), and:
\begin{itemize}
 \item \(f_{\mathbf C^k(\gamma;\vec\eta)}(\vec x)\simeq
   f_\gamma(f_{\eta_1}(\vec x),\ldots,f_{\eta_m}(\vec x))\), defined exactly
   when every \(f_{\eta_j}(\vec x)\) is defined and \(f_\gamma\) is defined at
   the resulting values;
 \item \(f_{\mathbf R(\gamma,\eta)}(\vec x,0)\simeq f_\gamma(\vec x)\) and
   \(f_{\mathbf R(\gamma,\eta)}(\vec x,n+1)\simeq
   f_\eta(\vec x,n,f_{\mathbf R(\gamma,\eta)}(\vec x,n))\), the latter defined
   exactly when \(f_{\mathbf R(\gamma,\eta)}(\vec x,n)\) is defined and
   \(f_\eta\) is defined there;
 \item \(f_{\mathbf M(\gamma)}(\vec x)\) is the least \(y\) such that
   \(f_\gamma(\vec x,y)\simeq0\) and \(f_\gamma(\vec x,z)\) is defined and
   nonzero for every \(z<y\), and is undefined when there is no such \(y\).
\end{itemize}
The \emph{graph} of \(\rho\) is
\(\mathrm{Gr}_\rho=\{(\vec x,y)\in\N^{\mathrm{ar}(\rho)+1}:f_\rho(\vec x)\simeq y\}\).
\index{recursion term}
\end{hbdefinition}

The operations on partial maps in this definition are those of
\cref{def:p04-partial-recursive}, and on total maps they are those of
\cref{sec:turing-machines-02}.  Terms admit arity \(0\) so that the zero
function \(Z_0\) and primitive recursion with \(k=0\) of
\cref{sec:turing-machines-02} have terms.

\begin{hblemma}{thm:p02-w02-p02-041-coverage}{proved=annex:p02-w02-p02-041-coverage}{Every recursive function has a term}
\begin{enumerate}
 \item[(a)] For every \(k\ge1\) and every partial recursive
   \(f:\N^k\rightharpoonup\N\) in the sense of
   \cref{def:p04-partial-recursive}, some term \(\rho\) of arity \(k\)
   satisfies \(f_\rho=f\).
 \item[(b)] For every \(k\in\N\) and every primitive recursive
   \(f:\N^k\to\N\) in the sense of \cref{sec:turing-machines-02}, some term
   \(\rho\) of arity \(k\) satisfies \(f_\rho=f\).
\end{enumerate}
\end{hblemma}

\subsection{Second-order definitions of graphs}
\label{sec:p02-w02-p02-041-graphs}

\begin{hbdefinition}{def:p02-w02-p02-041-graph-sentence}{Graph conditions and the divergence sentence}
Write \(x_i\) for the individual variable with code \(N_0(i)\) and
\(X^n_i\) for the relation variable with code \(N_{15}(n,i)\)
(\cref{def:p02-w02-p02-040-so-coding}); \(\varphi\leftrightarrow\psi\)
abbreviates \((\varphi\to\psi)\land(\psi\to\varphi)\), and
conjunctions are grouped to the right.  For a term \(\sigma\) of arity \(k\)
put \(X_\sigma\defeq X^{k+1}_{\ulcorner\sigma\urcorner}\), write \(\vec x\) for
\(x_1,\ldots,x_k\) and \(\forall\vec x\) for
\(\forall x_1\cdots\forall x_k\).  The \emph{graph condition}
\(\mathsf{Def}_\sigma\) is the following formula.
\begin{enumerate}
 \item \(\sigma=\mathbf Z^k\):\quad
   \(\forall\vec x\,\forall x_0\,(X_\sigma(\vec x,x_0)\leftrightarrow x_0=0)\).
 \item \(\sigma=\mathbf S\):\quad
   \(\forall x_1\,\forall x_0\,(X_\sigma(x_1,x_0)\leftrightarrow x_0=S x_1)\).
 \item \(\sigma=\mathbf P^k_i\):\quad
   \(\forall\vec x\,\forall x_0\,(X_\sigma(\vec x,x_0)\leftrightarrow x_0=x_i)\).
 \item \(\sigma=\mathbf C^k(\gamma;\eta_1,\ldots,\eta_m)\):
   \[
   \forall\vec x\,\forall x_0\,\Bigl(X_\sigma(\vec x,x_0)\leftrightarrow
   \exists x_{k+1}\cdots\exists x_{k+m}\bigl(X_{\eta_1}(\vec x,x_{k+1})\land\cdots\land
   X_{\eta_m}(\vec x,x_{k+m})\land X_\gamma(x_{k+1},\ldots,x_{k+m},x_0)\bigr)\Bigr).
   \]
 \item \(\sigma=\mathbf R(\gamma,\eta)\) with \(\mathrm{ar}(\gamma)=k\): the
   conjunction of
   \[
   \forall\vec x\,\forall x_0\,\bigl(X_\sigma(\vec x,0,x_0)\leftrightarrow X_\gamma(\vec x,x_0)\bigr)
   \]
   and
   \[
   \forall\vec x\,\forall x_{k+1}\,\forall x_0\,\Bigl(X_\sigma(\vec x,S x_{k+1},x_0)\leftrightarrow
   \exists x_{k+2}\bigl(X_\sigma(\vec x,x_{k+1},x_{k+2})\land
   X_\eta(\vec x,x_{k+1},x_{k+2},x_0)\bigr)\Bigr).
   \]
 \item \(\sigma=\mathbf M(\gamma)\) with \(\mathrm{ar}(\gamma)=k+1\):
   \[
   \forall\vec x\,\forall x_0\,\Bigl(X_\sigma(\vec x,x_0)\leftrightarrow
   \bigl(X_\gamma(\vec x,x_0,0)\land\forall x_{k+1}\bigl(\exists x_{k+2}\,(x_{k+1}+S x_{k+2}=x_0)
   \to\exists x_{k+3}\,(X_\gamma(\vec x,x_{k+1},x_{k+3})\land\neg\,x_{k+3}=0)\bigr)\bigr)\Bigr).
   \]
\end{enumerate}
The relation variables free in \(\mathsf{Def}_\sigma\) are \(X_\sigma\) and
\(X_\tau\) for the immediate subterms \(\tau\) of \(\sigma\).  The
\emph{subterm list} \(\mathrm{Sub}(\rho)\) is \([\rho]\) for clauses 1--3,
and otherwise \([\rho]\) followed by the subterm lists of the immediate
subterms in the order displayed (repetitions kept).  For a term \(\rho\) of
arity \(1\) with \(\mathrm{Sub}(\rho)=[\sigma_1,\ldots,\sigma_n]\), the
\emph{divergence sentence} is
\[
 \Theta_\rho\defeq
 \mathsf{PA}_2\to\forall X_{\sigma_1}\cdots\forall X_{\sigma_n}
 \Bigl(\bigl(\mathsf{Def}_{\sigma_1}\land\cdots\land\mathsf{Def}_{\sigma_n}\bigr)
 \to\neg\exists x_0\,X_\rho\bigl(\overline{\ulcorner\rho\urcorner},x_0\bigr)\Bigr),
\]
where \(\overline{\ulcorner\rho\urcorner}\) is the numeral of the code of
\(\rho\).  It is a second-order sentence of the signature
\(\{0,S,+,\times\}\).
\index{graph condition}
\index{divergence sentence}
\end{hbdefinition}

Distinct subterms have distinct codes, hence distinct relation variables; a
subterm occurring twice receives the same variable and the same condition, so
repetitions in \(\mathrm{Sub}(\rho)\) change neither the meaning nor the
proofs below.

\begin{hblemma}{thm:p02-w02-p02-041-iso}{proved=annex:p02-w02-p02-041-iso}{Isomorphism invariance under full semantics}
Let \(f\) be an isomorphism from a structure \(\mathcal A\) onto a structure
\(\mathcal B\) of the same signature.  For a full second-order assignment
\(s\) over \(\mathcal A\) let \(f\circ s\) send each individual variable
\(x\) to \(f(s(x))\) and each \(n\)-ary relation variable \(X\) to
\(\{(f(a_1),\ldots,f(a_n)):(a_1,\ldots,a_n)\in s(X)\}\).  Then for every
second-order formula \(\varphi\),
\[
 \mathcal A,s\entails\varphi\quad\Longleftrightarrow\quad\mathcal B,f\circ s\entails\varphi .
\]
\end{hblemma}

\begin{hblemma}{thm:p02-w02-p02-041-graphs}{proved=annex:p02-w02-p02-041-graphs}{The graph conditions have exactly one solution}
Let \(\mathcal N=(\N,0,n\mapsto n+1,+,\times)\), let \(\rho\) be a term and
let \(s\) be a full second-order assignment over \(\mathcal N\).
\begin{enumerate}
 \item If \(s(X_\tau)=\mathrm{Gr}_\tau\) for every immediate subterm \(\tau\)
   of a term \(\sigma\), then
   \(\mathcal N,s\entails\mathsf{Def}_\sigma\) iff
   \(s(X_\sigma)=\mathrm{Gr}_\sigma\).
 \item \(\mathcal N,s\entails\mathsf{Def}_\sigma\) for every \(\sigma\) in
   \(\mathrm{Sub}(\rho)\) iff \(s(X_\sigma)=\mathrm{Gr}_\sigma\) for every
   \(\sigma\) in \(\mathrm{Sub}(\rho)\).
\end{enumerate}
\end{hblemma}

\begin{hbtheorem}{thm:p02-w02-p02-041-reduction}{proved=annex:p02-w02-p02-041-reduction}{The divergence sentence is valid exactly at divergence}
For every term \(\rho\) of arity \(1\),
\[
 \entails_{\mathrm{full}}\Theta_\rho
 \quad\Longleftrightarrow\quad
 f_\rho(\ulcorner\rho\urcorner)\text{ is undefined}.
\]
\end{hbtheorem}

\subsection{Effectiveness and the diagonal argument}
\label{sec:p02-w02-p02-041-diagonal}

\begin{hblemma}{thm:p02-w02-p02-041-effective}{proved=annex:p02-w02-p02-041-effective}{The divergence sentence is primitive recursive in the term}
Define \(\mathrm{rd}:\N\to\N\) by
\(\mathrm{rd}(c)=\ulcorner\Theta_\rho\urcorner\) when \(c=\ulcorner\rho\urcorner\)
for a term \(\rho\) of arity \(1\), and \(\mathrm{rd}(c)=N_6()\), the code of
\(\bottom\), otherwise.  Then \(\mathrm{rd}\) is primitive recursive.  So
are the test \(\mathrm{RT}(c)\), which holds iff \(c\) is the code of a
term, and the arity function \(\mathrm{ar}(c)\), which returns the arity of the
coded term and \(0\) on other inputs.
\end{hblemma}

\begin{hbtheorem}{thm:p02-w02-p02-041-no-recursive-domain}{proved=annex:p02-w02-p02-041-no-recursive-domain}{No recursive enumeration of full second-order validity}
Let \(\mathrm{Val}_2\) be the set of codes of full second-order validities over
\(\{0,S,+,\times\}\), as in \cref{def:p02-w02-p02-040-so-coding}.
\begin{enumerate}
 \item[(a)] No term \(\beta\) of arity \(1\) satisfies
   \(\operatorname{dom}(f_\beta)=\mathrm{Val}_2\).  Hence
   \(\mathrm{Val}_2\) is not the domain of any partial recursive function.
 \item[(b)] Let \(\mathsf D\) be a calculus whose set of theorem codes is the
   domain of a partial recursive function and which satisfies
   \(\mathsf D\derives\sigma\Leftrightarrow\entails_{\mathrm{full}}\sigma\) for every
   second-order sentence \(\sigma\) of this signature.  No such \(\mathsf D\)
   exists.
 \item[(c)] \(\mathrm{Val}_2\) is not computably enumerable.
 \item[(d)] No calculus whose set of theorem codes is computably enumerable is
   sound and complete for full validity in the sense of (b).
\end{enumerate}
\end{hbtheorem}

Clauses (c) and (d) are clauses 4 and 5 of
\cref{thm:p02-w02-p02-040-no-effective-calculus}.  Clause (c) follows from (a)
because a computably enumerable set is the domain of a partial computable
function, and \cref{thm:p04-tm-partial-recursive} identifies partial
computable with partial recursive functions; clause (d) follows from (b) in the
same way.  That identification is the one imported step: the handbook cites
it to Kleene's coextensiveness theorem for partial recursive and computable
functions \cite[\S69, Theorem~XXX, p.~376]{Kleene1952}, which is stated for
Kleene's own machine formulation.  The reduction runs
from the divergence problem of recursion terms to full validity; true
arithmetic does not occur in it.

\begin{hbexample}{def:p02-w02-p02-041-example}{A divergent term and a convergent one}
Let \(\gamma_0=\mathbf C^2(\mathbf S;\mathbf P^2_2)\), a term of arity
\(2\) with \(f_{\gamma_0}(x,y)=y+1\), and let \(\rho_0=\mathbf M(\gamma_0)\),
of arity \(1\).  Then
\(\mathrm{Sub}(\rho_0)=[\rho_0,\gamma_0,\mathbf S,\mathbf P^2_2]\), and the four
conditions are
\[
\begin{aligned}
 \mathsf{Def}_{\mathbf S}&:\ \forall x_1\forall x_0\,(X_{\mathbf S}(x_1,x_0)\leftrightarrow x_0=Sx_1),\\
 \mathsf{Def}_{\mathbf P^2_2}&:\ \forall x_1\forall x_2\forall x_0\,(X_{\mathbf P^2_2}(x_1,x_2,x_0)\leftrightarrow x_0=x_2),\\
 \mathsf{Def}_{\gamma_0}&:\ \forall x_1\forall x_2\forall x_0\,\bigl(X_{\gamma_0}(x_1,x_2,x_0)\leftrightarrow
   \exists x_3\,(X_{\mathbf P^2_2}(x_1,x_2,x_3)\land X_{\mathbf S}(x_3,x_0))\bigr),\\
 \mathsf{Def}_{\rho_0}&:\ \forall x_1\forall x_0\,\Bigl(X_{\rho_0}(x_1,x_0)\leftrightarrow
   \bigl(X_{\gamma_0}(x_1,x_0,0)\land\forall x_2\bigl(\exists x_3\,(x_2+Sx_3=x_0)\to
   \exists x_4\,(X_{\gamma_0}(x_1,x_2,x_4)\land\neg\,x_4=0)\bigr)\bigr)\Bigr).
\end{aligned}
\]
In \(\mathcal N\) the first three conditions force
\(X_{\gamma_0}=\{(x,y,y+1):x,y\in\N\}\), which contains no triple with last
entry \(0\); the fourth then forces \(X_{\rho_0}=\emptyset\).  So
\(f_{\rho_0}(\ulcorner\rho_0\urcorner)\) is undefined, and by
\cref{thm:p02-w02-p02-041-reduction} the sentence \(\Theta_{\rho_0}\) is a full
validity.

For the term \(\mathbf S\) of arity \(1\), \(f_{\mathbf S}(\ulcorner\mathbf S\urcorner)=\ulcorner\mathbf S\urcorner+1\)
is defined.  The structure \(\mathcal N\) with \(X_{\mathbf S}\) assigned the
graph of the successor satisfies \(\mathsf{PA}_2\) and
\(\mathsf{Def}_{\mathbf S}\) but not
\(\neg\exists x_0\,X_{\mathbf S}(\overline{\ulcorner\mathbf S\urcorner},x_0)\), so
\(\Theta_{\mathbf S}\) is not valid.
\end{hbexample}

\begin{hbwarning}{def:p02-w02-p02-041-limits}{What the argument uses and what it does not give}
Full semantics is used twice: in \cref{thm:p02-w02-p02-040-pa2-categorical}, to
replace an arbitrary model of \(\mathsf{PA}_2\) by \(\mathcal N\), and in
\cref{thm:p02-w02-p02-041-graphs}, where the graph \(\mathrm{Gr}_\sigma\) must
be an admissible value of \(X_\sigma\).  Under general (Henkin) semantics
neither step is available, and the theorem does not apply: the matched
calculus of \cref{def:p02-sol-derived-calculus} is complete for that
semantics (\cref{thm:p02-sol-derived-adequacy}).  The theorem concerns the
fixed signature \(\{0,S,+,\times\}\) and the coding of
\cref{def:p02-w02-p02-040-so-coding}.  It does not show that true
arithmetic is not computably enumerable (clause~3 of
\cref{thm:p02-w02-p02-040-no-effective-calculus}); that clause still rests on
\cref{thm:p04-ce-sigma-one}.
\end{hbwarning}

\section{Truth preservation through incidence encoding and extensional decoding}
\label{sec:p02-w02-p02-041-decoding}

\emph{Purpose.}  The matched relational calculus of
\cref{def:p02-sol-derived-calculus} is defined by translation into
many-sorted first-order logic, and its strong adequacy
(\cref{thm:p02-sol-derived-adequacy}) moves models in both directions: a
Henkin structure is encoded as a many-sorted structure in which relations are
objects of the sorts \(S_n\), and a many-sorted structure is decoded into a
Henkin structure by replacing each relation-sort object with the set of tuples
incident to it.  This section states, as separate results, the two
truth-preservation invariants that carry formulas across these maps, together
with the coincidence lemma for Henkin semantics on which the schema clauses
depend, and assembles them into both directions of the adequacy theorem.  All
cases of each induction are written out in the annex.

\emph{Inputs.}  The relational second-order language
(\cref{sec:second-order-logic-01}) with Henkin structures and matching
assignments (\cref{sec:second-order-logic-05}); the incidence signature, the
translation \(\varphi\mapsto\varphi^\sharp\) and the tagged structure
\(\mathcal H^\sharp\) of \cref{def:p02-sol-incidence-translation}; strict
many-sorted structures (\cref{def:p02-many-sorted-structure}); the matched
model class of \cref{def:p02-sol-matched-henkin}; first-order term values and
satisfaction clauses (\cref{sec:fol-semantics-01,sec:fol-semantics-02}).
\emph{Conclusions.}
\Cref{thm:p02-w02-p02-041-henkin-coincidence,thm:p02-w02-p02-041-tagging,thm:p02-w02-p02-041-collapse,thm:p02-w02-p02-041-adequacy-route}.
\emph{Standing.}  The constructions are those already used in the proofs of
\cref{thm:p02-sol-incidence-satisfaction,thm:p02-sol-derived-adequacy}; the
results below supply every case of the inductions those proofs invoke.

\emph{Conventions.}  Object terms are first-order terms of the signature, and
their value \(\operatorname{val}^{\mathcal H}_s(t)\) in a Henkin structure is
the first-order value in its object reduct under the restriction of \(s\) to
object variables.  The connectives are
\(\bottom,\neg,\land,\lor,\to,\leftrightarrow\), each with the
truth-functional clause of \cref{sec:fol-semantics-02} (for
\(\leftrightarrow\): true iff both sides have the same truth value) in every
semantics used here.  The incidence translation uses each source object
variable as a variable of sort \(D\) and each source variable \(X^n\) as a
variable of sort \(S_n\); a \emph{sorted assignment} of a many-sorted
structure sends the variables of each sort into that sort's domain.  In
\(\mathcal H^\sharp\) the object constants and functions are transported along
the object tag:
\(c^{\mathcal H^\sharp}=(D,c^{\mathcal H})\) and
\(f^{\mathcal H^\sharp}((D,a_1),\ldots,(D,a_k))=(D,f^{\mathcal H}(a_1,\ldots,a_k))\).
A Henkin structure \emph{satisfies} a formula when it does so under every
matching assignment.

\begin{hbdefinition}{def:p02-w02-p02-041-collapse}{Extensional collapse of an incidence structure}
Let \(\mathcal A\) be a strict many-sorted structure for the incidence
signature of \cref{def:p02-sol-incidence-translation}, with domains
\(D^{\mathcal A}\) and \(S_n^{\mathcal A}\) (\(n\geq1\)).  For
\(r\in S_n^{\mathcal A}\) put
\[
 \operatorname{ext}_n(r)\defeq
 \{(a_1,\ldots,a_n)\in(D^{\mathcal A})^n:
   (r,a_1,\ldots,a_n)\in E_n^{\mathcal A}\}.
\]
The \emph{extensional collapse} \(\mathcal H_{\mathrm{ext}}\) of
\(\mathcal A\) has object domain \(M\defeq D^{\mathcal A}\), relation domains
\[
 \mathcal R_n\defeq\operatorname{ext}_n[S_n^{\mathcal A}]
 =\{\operatorname{ext}_n(r):r\in S_n^{\mathcal A}\},
\]
object constants and functions \(c^{\mathcal H_{\mathrm{ext}}}=c^{\mathcal
A}\) and \(f^{\mathcal H_{\mathrm{ext}}}=f^{\mathcal A}\), and predicate
constants \(P^{\mathcal H_{\mathrm{ext}}}=\operatorname{ext}_n(\widehat
P^{\mathcal A})\).  For a sorted assignment \(s\) of \(\mathcal A\), the
\emph{decoded assignment} is given by
\(s^{\mathrm{ext}}(x)\defeq s(x)\) and
\(s^{\mathrm{ext}}(X^n)\defeq\operatorname{ext}_n(s(X^n))\).
The map \(\operatorname{ext}_n\) need not be injective: two relation-sort
objects with the same incident tuples collapse to one relation.
\index{extensional collapse}
\end{hbdefinition}

\begin{hblemma}{thm:p02-w02-p02-041-henkin-coincidence}{proved=annex:p02-w02-p02-041-henkin-coincidence}{Coincidence for Henkin semantics}
Let \(\mathcal H\) be a relational Henkin structure, \(\varphi\) a relational
second-order formula, and \(s,t\) matching assignments that agree on every
object and relation variable free in \(\varphi\).  Then
\(\mathcal H,s\entails\varphi\) iff \(\mathcal H,t\entails\varphi\).  In
particular a sentence true in \(\mathcal H\) under one matching assignment is
true under every matching assignment.
\end{hblemma}

\begin{hblemma}{thm:p02-w02-p02-041-tagging}{proved=annex:p02-w02-p02-041-tagging}{Tagging invariant for the incidence encoding}
Let \(\mathcal H=(M,(\mathcal R_n)_{n\geq1},\ldots)\) be a relational Henkin
structure with \(P^{\mathcal H}\in\mathcal R_n\) for every \(n\)-ary predicate
constant \(P\), and let \(s\) be a matching assignment.
\begin{enumerate}
\item[(a)] \(\mathcal H^\sharp\) is a strict many-sorted structure for the
incidence signature, and \(s^\sharp\) is a sorted assignment of it.
\item[(b)] For every object term \(t\),
\(\operatorname{val}^{\mathcal H^\sharp}_{s^\sharp}(t)=(D,\operatorname{val}^{\mathcal H}_s(t))\).
\item[(c)] For every relational second-order formula \(\varphi\),
\(\mathcal H,s\entails\varphi\) iff
\(\mathcal H^\sharp,s^\sharp\entails\varphi^\sharp\).
\item[(d)] Every sorted assignment of \(\mathcal H^\sharp\) is \(u^\sharp\) for
exactly one matching assignment \(u\) of \(\mathcal H\).
\item[(e)] \(\mathcal H\) is an \(\mathsf H_{\mathsf{CC}}\)-model iff
\(\mathcal H^\sharp\) satisfies, under every sorted assignment, every member of
\(T_{\mathsf{CC}}^\sharp\).
\end{enumerate}
\end{hblemma}

Clauses (c) and (e) together are \cref{thm:p02-sol-incidence-satisfaction};
the lemma adds the term clause and the assignment correspondence on which its
proof depends.

\begin{hbtheorem}{thm:p02-w02-p02-041-collapse}{proved=annex:p02-w02-p02-041-collapse}{Truth preservation under extensional collapse}
Let \(\mathcal A\) be a strict many-sorted structure for the incidence
signature, \(s\) a sorted assignment of \(\mathcal A\), and
\(\mathcal H_{\mathrm{ext}}\), \(s^{\mathrm{ext}}\) as in
\cref{def:p02-w02-p02-041-collapse}.
\begin{enumerate}
\item[(a)] \(\mathcal H_{\mathrm{ext}}\) is a relational Henkin structure,
\(P^{\mathcal H_{\mathrm{ext}}}\in\mathcal R_n\) for every \(n\)-ary predicate
constant \(P\), and \(s^{\mathrm{ext}}\) is a matching assignment.
\item[(b)] For every object term \(t\),
\(\operatorname{val}^{\mathcal H_{\mathrm{ext}}}_{s^{\mathrm{ext}}}(t)
=\operatorname{val}^{\mathcal A}_s(t)\).
\item[(c)] For every relational second-order formula \(\varphi\),
\[
 \mathcal A,s\entails\varphi^\sharp
 \quad\Longleftrightarrow\quad
 \mathcal H_{\mathrm{ext}},s^{\mathrm{ext}}\entails\varphi.
\]
\item[(d)] If \(\mathcal A,s\entails\theta\) for every
\(\theta\in T_{\mathsf{CC}}^\sharp\), then \(\mathcal H_{\mathrm{ext}}\) is an
\(\mathsf H_{\mathsf{CC}}\)-model.
\end{enumerate}
\end{hbtheorem}

The relation-quantifier case of (c) is where the collapse matters: a
quantifier of sort \(S_n\) in \(\mathcal A\) ranges over objects \(r\), while
the relation quantifier of \(\mathcal H_{\mathrm{ext}}\) ranges over their
extensions, and the two clauses agree because \(\mathcal R_n\) is exactly the
image \(\operatorname{ext}_n[S_n^{\mathcal A}]\).  Duplicate objects with one
extension change neither side, because no translated formula compares
relation-sort objects other than through \(E_n\).  The proof uses no choice
principle: each case instantiates one existential statement at a time.

\begin{hbcorollary}{thm:p02-w02-p02-041-adequacy-route}{proved=annex:p02-w02-p02-041-adequacy-route}{Matched adequacy through the encoding and the decoding}
Under the hypotheses of \cref{thm:p02-sol-derived-adequacy}, for every set
\(\Gamma\) of formulas and every formula \(\varphi\),
\[
 T_{\mathsf{CC}}^\sharp\cup\Gamma^\sharp\derives_{\mathsf{MSFOL}}\varphi^\sharp
 \quad\Longleftrightarrow\quad
 \Gamma\entails_{\mathrm H}^{\mathsf{CC}}\varphi .
\]
The forward direction uses \cref{thm:p02-w02-p02-041-tagging}; the backward
direction uses \cref{thm:p02-w02-p02-041-collapse}.
\end{hbcorollary}

By \cref{def:p02-sol-derived-calculus} the left side is
\(\Gamma\derives_{\mathsf{SOL}_{\mathrm H}^{\mathsf{CC}}}\varphi\), so this is
the displayed equivalence of \cref{thm:p02-sol-derived-adequacy}.  The
admissibility of the six sheet additions stated in that theorem is a separate
part of its annex proof and is not repeated here.

\begin{hbexample}{def:p02-w02-p02-041-collapse-example}{A collapse that identifies two relation objects}
Take the signature with one object constant \(c\) and no function or predicate
constants, and let \(\mathcal A\) have \(D^{\mathcal A}=\{0,1\}\),
\(c^{\mathcal A}=0\), \(S_1^{\mathcal A}=\{r_1,r_2,r_3\}\) and
\(E_1^{\mathcal A}=\{(r_1,0),(r_2,0)\}\); for \(n\geq2\) let
\(S_n^{\mathcal A}=\{q_n\}\) with no incident tuple.  Then
\(\operatorname{ext}_1(r_1)=\operatorname{ext}_1(r_2)=\{0\}\) and
\(\operatorname{ext}_1(r_3)=\varnothing\), so
\(\mathcal R_1=\{\{0\},\varnothing\}\), and
\(\mathcal R_n=\{\varnothing\}\) for \(n\geq2\).

For \(\varphi=\exists X^1\,(X^1(c)\land\forall x\,(X^1(x)\to x=c))\) the
translation is
\[
 \varphi^\sharp=\exists X^1{:}S_1\,\bigl(E_1(X^1,c)\land\forall x{:}D\,(E_1(X^1,x)\to x=c)\bigr).
\]
In \(\mathcal A\) the object \(r_1\) witnesses \(\varphi^\sharp\); in
\(\mathcal H_{\mathrm{ext}}\) the relation \(\{0\}=\operatorname{ext}_1(r_1)\)
witnesses \(\varphi\), and both sides are true, as clause (c) requires.  For
\(\psi=\forall X^1\,(X^1(c)\lor\forall x\,\neg X^1(x))\) both sides are true
as well: \(r_1\) and \(r_2\) are incident with \(0\) and \(r_3\) with nothing,
while \(\{0\}\) and \(\varnothing\) are the only members of \(\mathcal R_1\).

The many-sorted sentence
\(\exists X{:}S_1\,\exists Y{:}S_1\,\bigl(\neg X=Y\land\forall x{:}D\,(E_1(X,x)\leftrightarrow E_1(Y,x))\bigr)\)
is true in \(\mathcal A\) (take \(r_1,r_2\)), whereas any two coextensive
members of \(\mathcal R_1\) are equal.  This sentence is not \(\chi^\sharp\)
for any source formula \(\chi\), since the translation never produces an
equation between relation-sort terms, so clause (c) says nothing about it.
\end{hbexample}

\begin{hbwarning}{def:p02-w02-p02-041-collapse-limits}{What the decoding does not preserve}
The invariant of \cref{thm:p02-w02-p02-041-collapse}(c) holds for formulas in
the image of the translation, not for arbitrary many-sorted formulas: as the
example shows, a sentence with relation-sort equality can change truth value.
The collapse \(\mathcal H_{\mathrm{ext}}\) encodes to
\(\mathcal H_{\mathrm{ext}}^\sharp\), which is isomorphic to \(\mathcal A\)
when every \(\operatorname{ext}_n\) is injective (by
\(r\mapsto(S_n,\operatorname{ext}_n(r))\) and \(a\mapsto(D,a)\)) and not
otherwise: in \(\mathcal H_{\mathrm{ext}}^\sharp\) distinct objects of sort
\(S_n\) have distinct incident tuples, while a non-injective
\(\operatorname{ext}_n\) gives two distinct coextensive objects in
\(\mathcal A\), and an isomorphism preserves that difference.  The decoded structure has the selected relation domains
\(\operatorname{ext}_n[S_n^{\mathcal A}]\), in general a proper subset of
\(\mathcal P(M^n)\); the decoding yields Henkin countermodels and says nothing
about full validity, for which no calculus with a recursively enumerable set of
theorems is sound and complete
(\cref{thm:p02-w02-p02-041-no-recursive-domain}).
\end{hbwarning}

\section{Church's formulation: notation, qualifications, formulas 30}
\label{sec:p02-w02-p02-041-church-supplement}

\emph{Purpose.}  This section completes the presentation of Church's system in
\cref{sec:p02-w02-p02-040-church} on four points: the worked examples of the
notational conventions, two qualifications the source attaches to its
vocabulary, the exact attribution of the non-definability remark about the
same-type predecessor, and a proof that not every formula 30\(^\alpha\) is a
theorem.  \emph{Inputs.}  \Cref{def:p02-w02-p02-040-church-types,def:p02-w02-p02-040-church-abbrev,def:p02-w02-p02-040-church-rules,def:p02-w02-p02-040-church-axioms,def:p02-w02-p02-040-church-standard}
and \cref{thm:p02-w02-p02-040-church-soundness}.  \emph{Conclusions.}
\Cref{def:p02-w02-p02-041-church-parentheses,def:p02-w02-p02-041-church-brackets,thm:p02-w02-p02-041-church-30}.
\emph{Standing.}  Examples and qualifications are the source's, restated; the
non-theoremhood proof is the handbook's.

\emph{Currying.}  The source credits the treatment of functions of several
variables as functions of one variable whose values are again functions to
Sch\"onfinkel \cite[\S1, p.~57, n.~6]{Church1940}; this is the convention by
which \(o\iota\iota\) is the type of two-place propositional functions of
individuals in \cref{def:p02-w02-p02-040-church-types}.

\begin{hbexample}{def:p02-w02-p02-041-church-parentheses}{Omitting parentheses and subscripts}
Parentheses may be dropped when no ambiguity results; they are restored so
that the result is well-formed and, apart from that, by association to the
left.  Type subscripts may be shortened as in
\cref{def:p02-w02-p02-040-church-types}, and a subscript \(o\) may be left
out entirely, so that a small italic letter with no subscript is a variable of
type \(o\).  The source's two examples are \cite[\S2, p.~58]{Church1940}:
\[
\begin{aligned}
 a_{\iota'}b_{\iota\iota}(c_{\iota\iota}d_\iota)
 &\quad\text{abbreviates}\quad
 ((a_{((\iota\iota)(\iota\iota))}b_{(\iota\iota)})(c_{\iota\iota}d_\iota)),\\
 \lambda b_{\iota\iota}\lambda c_{\iota\iota}(a_{\iota'}b_{\iota\iota}(c_{\iota\iota}d_\iota))
 &\quad\text{abbreviates}\quad
 (\lambda b_{\iota\iota}(\lambda c_{\iota\iota}((a_{((\iota\iota)(\iota\iota))}b_{(\iota\iota)})(c_{\iota\iota}d_\iota)))).
\end{aligned}
\]
In the first, \(a_{\iota'}\) has type \(\iota'=((\iota\iota)(\iota\iota))\), so
it is applied first to \(b_{\iota\iota}\), giving a formula of type
\(\iota\iota\), which is then applied to \((c_{\iota\iota}d_\iota)\), of type
\(\iota\); left association produces exactly this grouping.  In the second the
two abstraction operators are restored around the whole remaining formula.
\end{hbexample}

\begin{hbexample}{def:p02-w02-p02-041-church-brackets}{Restoring square brackets by rank}
The source illustrates the dot convention and the bracket ranks of
\cref{def:p02-w02-p02-040-church-abbrev} with
\cite[\S2, p.~59]{Church1940}
\[
 \sim p\supset q\supset.\,pq\vee rs\supset\sim.\,q\vee s\supset\sim p\sim r,
\]
which abbreviates
\[
 [[[\sim p]\supset q]\supset[[[pq]\vee[rs]]\supset[\sim[[q\vee s]\supset[[\sim p][\sim r]]]]]],
\]
which in turn abbreviates the well-formed formula
\[
\begin{aligned}
 &((A_{ooo}(N_{oo}((A_{ooo}(N_{oo}(N_{oo}p_o)))q_o)))((A_{ooo}(N_{oo}((A_{ooo}(N_{oo}((A_{ooo}(N_{oo}p_o))(N_{oo}q_o))))\\
 &\qquad(N_{oo}((A_{ooo}(N_{oo}r_o))(N_{oo}s_o))))))\\
 &\qquad(N_{oo}((A_{ooo}(N_{oo}((A_{ooo}q_o)s_o)))\\
 &\qquad\qquad(N_{oo}((A_{ooo}(N_{oo}(N_{oo}p_o)))(N_{oo}(N_{oo}r_o)))))))).
\end{aligned}
\]
\emph{First stage.}  Each dot opens a bracket whose scope runs as far right as
well-formedness allows, here to the end: the first dot encloses
\(pq\vee rs\supset\sim.\,q\vee s\supset\sim p\sim r\), and the dot after
\(\sim\) encloses \(q\vee s\supset\sim p\sim r\).  Within each part the
lowest-rank brackets are restored first, giving \([\sim p]\), \([\sim r]\),
\([pq]\), \([rs]\) and the conjunction \([[\sim p][\sim r]]\); then the
brackets of \(\vee\), giving \([q\vee s]\) and \([[pq]\vee[rs]]\); and last
those of \(\supset\), associating to the left, so that
\(\sim p\supset q\supset X\) becomes \([[[\sim p]\supset q]\supset X]\).
\emph{Second stage.}  Each bracketed abbreviation is replaced by its definiens:
\([\mathbf A\supset\mathbf B]\) by \([[\sim\mathbf A]\vee\mathbf B]\),
\([\mathbf A\mathbf B]\) by \([\sim[[\sim\mathbf A]\vee[\sim\mathbf B]]]\),
\([\mathbf A\vee\mathbf B]\) by \(A_{ooo}\mathbf A\mathbf B\) and
\([\sim\mathbf A]\) by \(N_{oo}\mathbf A\), after which every parenthesis of
application is written and every omitted subscript \(o\) restored.  For
instance \([[\sim p]\supset q]\) becomes \([[\sim[\sim p]]\vee q]\) and then
\(((A_{ooo}(N_{oo}(N_{oo}p_o)))q_o)\), which occurs, under \(N_{oo}\), at the start of the display above.
\end{hbexample}

\emph{Two qualifications on vocabulary.}  The source separates \emph{formal
theorems}, the theorems of the system, from \emph{syntactical theorems},
statements proved about the system in the syntax language, and it says that
it uses the word ``theorem'' ambiguously, for a proposition and also for a
sentence or formula expressing that proposition in some language
\cite[\S4, p.~61]{Church1940}.  The deduction theorem VII
(\cref{thm:p02-w02-p02-040-church-deduction}) is a syntactical theorem in this
sense.  The source also marks its sign \(\derives\) for provability on
assumptions as a syntactical notation that must not be confused with the
assertion sign used, with a different meaning, by Russell and earlier by Frege
\cite[\S5, p.~62]{Church1940}.

\begin{hbwarning}{def:p02-w02-p02-041-church-predecessor}{Attributed standing of the predecessor remark}
The source says that the same-type predecessor \(P_{\alpha'\alpha'}\) cannot
be defined without descriptions, that is without the selection constant, and
derives this from the independence of the axiom of infinity of its \S4
\cite[\S2, p.~59]{Church1940}.  In \S4 the axiom of infinity is not a single
formula: Axioms~7 and~8 taken together have that effect
\cite[\S4, p.~61]{Church1940}.  The remark therefore rests on the independence
of Axioms~7 and~8 jointly, not on Axiom~8 alone.  The source states
the non-definability claim without proof, and the handbook does not prove it;
it is a reported source claim.  What is proved is the pair of independence
results the source states on p.~61: clause~3 of
\cref{thm:p02-w02-p02-040-church-soundness} shows that Axiom~7 is not a
theorem of the system whose axioms are all the others, and likewise for
Axiom~8.  The passage from those results to the non-definability of
\(P_{\alpha'\alpha'}\) is the step the source leaves unproved.
\end{hbwarning}

The source also states that, as with 26\(^\alpha\), not all of the formulas
\[
 30^\alpha.\quad N_{o\alpha''}m_{\alpha''}\supset.\,N_{o\alpha''}n_{\alpha''}\supset.\,
 m_{\alpha''}S_{\alpha'\alpha'}0_{\alpha'}=n_{\alpha''}S_{\alpha'\alpha'}0_{\alpha'}
 \supset m_{\alpha''}=n_{\alpha''}
\]
are theorems \cite[\S6, p.~66]{Church1940}.  The following result proves this
for the type \(o\).

\begin{hbtheorem}{thm:p02-w02-p02-041-church-30}{proved=annex:p02-w02-p02-041-church-30}{Not every formula 30 is a theorem}
Work in ZFC.  The formula 30\(^o\) is not valid in the standard model
\(\mathcal S(\N,c)\) of \cref{def:p02-w02-p02-040-church-standard}.  Hence
30\(^o\) is not a theorem of Church's system, and not all formulas
30\(^\alpha\) are theorems.
\end{hbtheorem}

In \(\mathcal S(\N,c)\) the numerals \(1_{o'}\) and \(3_{o'}\) denote the same
element, because every function from a two-element set to itself equals its
own third iterate, while \(1_{o''}\) and \(3_{o''}\) denote different elements;
30\(^o\) fails at \(m_{o''}=1_{o''}\), \(n_{o''}=3_{o''}\).  This is the same
finite-type mechanism that defeats 26\(^\alpha\) in
\cref{thm:p02-w02-p02-040-church-soundness}(4), applied one type level higher.
